% !TEX root = ../../main_power.tex
% ---------------------------------------------------------------------
%  Formelsammlung und Argumentationshilfe
% ---------------------------------------------------------------------
\section{Referenz}
\label{sec:ref}

\subsection{Formeln auf einen Blick}
\begin{formel}[Alles, was gerechnet wird]
\begin{align*}
P &= U \cdot I \\
E &= \int_0^T P(t)\,dt \;=\; \int_0^T I(t)\,V(t)\,dt \\
I &= \frac{dQ}{dt} \;=\; C \cdot \frac{dV}{dt} \\
P_{total} &= \alpha C_{load} V_{dd}^{2} f \;+\; t_{SC} V_{dd} I_{peak}
   \;+\; V_{dd} I_{leak}
\end{align*}
Mittelwerte: $E = P_{avg} \cdot T$, $P_{avg} = U \cdot I_{avg}$.
Einheiten: 1\,Wh $=$ 3600\,Ws; 1\,\textmu Wh $=$ 3.6\,mWs.
\end{formel}

\begin{bsp}[Typische Rechnung: Batterielaufzeit]
Ein Gerät wacht jede Sekunde für 100\,ms auf (5\,mA) und schläft 900\,ms
(10\,\textmu A), bei 3.0\,V.\\
$I_{avg} = 0.1\cdot 5\,\text{mA} + 0.9\cdot 0.01\,\text{mA}
 = 0.5 + 0.009 = 0.509\,\text{mA}$\\
$P_{avg} = 3.0\,\text{V} \cdot 0.509\,\text{mA} = 1.53\,\text{mW}$\\
Mit 220\,mAh: $t = 220/0.509 \approx 432\,$h $\approx$ 18 Tage.\\[2pt]
\merkzeile Die Wachphase macht hier 98\,\% des Verbrauchs aus. Den
Schlafstrom weiter zu senken bringt fast nichts -- die Wachzeit zu
halbieren halbiert die Laufzeit\emph{kosten}. Immer zuerst ausrechnen,
welcher Summand dominiert.
\end{bsp}

\subsection{Rechnen mit Sleep-Modi}
\begin{formel}[Lohnt sich das Schlafen?]
Schlafen lohnt nur, wenn die Schlafzeit den Aufwach-Overhead trägt:
\[ t_{sleep} > t_{wake} \cdot \frac{P_{wake} - P_{sleep}}{P_{run} - P_{sleep}} \]
Praktisch: \T{die Energie zum Aufwachen muss kleiner sein als die
Energie, die man durch das Schlafen spart}. Bei einem 32-kHz-Quarz mit
einigen hundert Millisekunden Anschwingzeit ist ein Schlaf von 10\,ms
ein Verlustgeschäft.
\end{formel}
\begin{bsp}[Gear Shifting durchgerechnet]
Eine Aufgabe braucht $N$ Takte. Bei 20\,MHz dauert sie
$T_1 = N/20\,$MHz, bei 10\,MHz doppelt so lang.\\
\T{Dynamisch:} $P \propto f$ \ra halbe Leistung, doppelte Zeit
\ra \emph{gleiche Energie}.\\
\T{Statisch:} $P = V_{dd}I_{leak}$ ist von $f$ \emph{unabhängig}
\ra doppelte Zeit $=$ \textbf{doppelte} statische Energie.\\[2pt]
\merkzeile Den Takt allein zu senken bringt energetisch \emph{nichts}
und kann schaden. Es lohnt sich nur, wenn man gleichzeitig $V_{dd}$
senken kann (dann wirkt das Quadrat) oder wenn die Alternative Warten
statt Schlafen wäre. Das ist die wichtigste Fehlannahme im ganzen Thema.
\end{bsp}

\begin{rezept}[Eine Messung, die als Argument taugt]
\begin{enumerate}
\item \T{Ziel definieren} -- was genau soll kleiner werden, Leistung
      oder Energie, und in welchem Betriebszustand?
\item \T{Messpunkt festlegen} -- vor oder nach dem Regler? Nur der
      MCU oder das ganze Board? Der Wirkungsgrad des Reglers
      (beim tinyK22 80\,\% für die SD-Karte) steckt sonst
      unbemerkt in jeder Zahl.
\item \T{Basismessung} vor jeder Änderung, mit Zeitstempel und
      Screenshot ins Journal.
\item \T{Eine Variable pro Durchgang} ändern -- $V_{dd}$, dann
      Frequenz, dann Modus. Sonst lässt sich die Wirkung nicht
      zuordnen.
\item \T{min/avg/max und Dauer} notieren, nicht nur den Mittelwert.
      Bei Lastwechseln gehört der Verlauf dazu.
\item \T{Gegenrechnen} -- passt die Messung zur Theorie? Wenn
      $V_{dd}$ halbiert wird und der Strom nicht um etwa den Faktor 4
      fällt, ist entweder der statische Anteil dominant oder der
      Messaufbau falsch.
\end{enumerate}
\end{rezept}

\begin{falle}[Die drei häufigsten Messfehler]
\T{(a) Falscher Messbereich.} \textmu A im Schlaf und mA im Betrieb
passen nicht in denselben Bereich -- der Mittelwert wird dann von
Quantisierung dominiert.\\
\T{(b) Der Debugger speist mit.} Solange die Debug-Probe am Target
hängt, fliesst Strom über die SWD-Leitungen. Für die Endmessung
trennen oder den Pfad kennen.\\
\T{(c) Das Board ist nicht der Chip.} Regler, Debug-Chip, LEDs und
Pull-ups laufen weiter. Die gemessenen 3.39\,mA im \R{STOP}-Modus
(\ref{sec:messen}) sind überwiegend \emph{nicht} der K22.
\end{falle}

\subsection{Von der Aufgabe zur Massnahme}
\begin{ztab}{@{}lX@{}}
„Verbrauch halbieren“ & zuerst $V_{dd}$ (quadratisch), dann Wachzeit,
  dann Takt \ref{form:ptotal} \\
„warum wird es schlechter, obwohl ich den Takt senke“ & statischer
  Anteil bleibt, Arbeit dauert länger \ra mehr Zeit mal $V_{dd}I_{leak}$
  \ref{sec:statisch} \\
„Peripherie aus, aber es ändert nichts“ & Clock Gating statt Power
  Gating, oder parasitäre Speisung über Signalleitungen
  \ref{rez:gating} \\
„HardFault nach dem Sparen“ & Zugriff auf Peripherie ohne Takt
  \ref{rez:gearshift} \\
„UART spinnt nach Gear Shift“ & Prescaler nicht nachgezogen
  \ref{rez:gearshift} \\
„welcher Modus für diese Weckquelle“ & synchron $\to$ Sleep (NVIC),
  asynchron $\to$ Deepsleep (AWIC/LLWU) \ref{rez:modi} \\
„RTOS schläft nicht richtig“ & Idle Hook gesetzt, aber Tick weckt
  \ra Tickless Idle \ref{rez:tickless} \\
„Messung unplausibel“ & Innenwiderstand, Dynamikbereich, parasitäre
  Pfade \ref{falle:pwr} \\
\end{ztab}

\subsection{Semesterleistung}
Dreierteams, tinyK22 $+$ LPC804, 5\,\% der MEP. Abgabe auf ILIAS:
Präsentation (10\,min $+$ 5\,min Befragung) und Laborjournal, beides PDF.
Bewertet werden Technik, Resultate, Kompetenz, Überzeugungskraft,
Reflexion und Präsentationstechnik.
\merkzeile Vorgehen: mit Blinky starten, analysieren, \emph{messen und
dokumentieren}, verbessern, recherchieren (Datenblätter, SDK-Beispiele),
Schlüsse ziehen. Vorher Ziele definieren: Was soll erreicht werden? Was
und wo wird gemessen? Welchen Einfluss haben Schaltung, $V_{dd}$,
Frequenz und Software?
